\documentclass[11pt]{article}

% --- Packages ---
\usepackage[margin=1in]{geometry}
\usepackage{fontenc}
\usepackage[utf8]{inputenc}
\usepackage{lmodern}
\usepackage{hyperref}
\usepackage{xcolor}
\usepackage{enumitem}
\usepackage{booktabs}
\usepackage{listings}
\usepackage{amsmath}
\usepackage{float}

% --- Hyperref Setup ---
\hypersetup{
    colorlinks=true,
    linkcolor=blue!50!black,
    urlcolor=blue!50!black
}

% --- Code Block Styles ---
\definecolor{codebg}{RGB}{247,248,250}
\lstdefinestyle{pythoncode}{
    backgroundcolor=\color{codebg},
    basicstyle=\ttfamily\small,
    frame=single,
    rulecolor=\color{black!10},
    breaklines=true,
    columns=fullflexible,
    language=Python,
    keywordstyle=\color{blue},
    commentstyle=\color{green!50!black}
}

% --- Custom Commands ---
\newcommand{\code}[1]{\texttt{\detokenize{#1}}}

\title{\textbf{Technical Report: Advanced Feature Extraction for Visual Place Recognition}}
\author{Expert Engineering Group --- Project 13}
\date{April 24, 2026}

\begin{document}

\maketitle

\section{Introduction}
Visual Place Recognition (VPR) serves as the foundational "where am I?" component for autonomous robots navigating in GPS-denied environments. In a campus setting, this task is complicated by extreme lighting shifts (Day-to-Night) and structural changes. This report details the top 5 state-of-the-art (SOTA) feature extractors as of 2026 and provides a structured implementation flow aligned with the existing system architecture (Benchmark vs. Live paths).

\section{Analysis of Top 5 Feature Extractors}

\subsection{1. SelaVPR++: Seamless Adaptation and Deep Hashing}
\textbf{Core Mechanism:} SelaVPR++ leverages a frozen foundation model (DINOv2) and adds Multi-scale Convolution (MultiConv) adapters. These adapters are added after the Multi-Head Attention (MHA) layers and in parallel to the MLP layers in the transformer blocks.
\begin{itemize}
    \item \textbf{High-Efficiency Paradigm:} It utilizes a two-stage retrieval. Stage 1 uses 512-dim binary features (generated via similarity-constrained deep hashing) for rapid initial retrieval. Stage 2 uses 4096-dim floating-point features for precision re-ranking.
    \item \textbf{Performance:} Reported to be 6000$\times$ faster than TransVPR with superior robustness against dynamic campus interference like trucks or pedestrians.
\end{itemize}

\subsection{2. VPRTempo: Spiking Neural Network (SNN) with Temporal Coding}
\textbf{Core Mechanism:} VPRTempo is a neuromorphic approach that trades biological realism for computational speed. It uses temporal encoding, where the timing of a single spike represents pixel intensity, rather than the frequency of spikes.
\begin{itemize}
    \item \textbf{Edge Suitability:} It achieves >50 Hz inference on a standard CPU. This allows the robot to perform VPR even when GPU resources are saturated by other tasks.
    \item \textbf{Training:} Can be trained within minutes using Spike-Timing Dependent Plasticity (STDP).
\end{itemize}

\subsection{3. SALAD: Sinkhorn Algorithm for Locally Aggregated Descriptors}
\textbf{Core Mechanism:} SALAD reformulates the aggregation of local features into global descriptors as an optimal transport problem solved via the Sinkhorn Algorithm.
\begin{itemize}
    \item \textbf{The "Dustbin" Cluster:} It introduces a specific cluster designed to discard non-informative features (like sky, pedestrians, or moved furniture). This ensures the global descriptor focuses exclusively on stable architectural landmarks.
    \item \textbf{Result:} Achieves unprecedented recall on MSLS and Nordland benchmarks.
\end{itemize}

\subsection{4. AnyLoc: Universal facet-based Retrieval}
\textbf{Core Mechanism:} AnyLoc identifies that specific "facets" (hooks from internal ViT layers) in pre-trained models are highly discriminative for VPR.
\begin{itemize}
    \item \textbf{Zero-Shot Capability:} It requires no local training or fine-tuning, making it the ideal baseline for immediate deployment on a new campus.
\end{itemize}

\subsection{5. MixVPR: MLP-based Global Descriptor Mixing}
\textbf{Core Mechanism:} MixVPR treats feature maps as sequences and uses simple MLP layers to mix spatial and channel information, avoiding the overhead of Attention mechanisms.
\begin{itemize}
    \item \textbf{Memory Efficiency:} It produces highly discriminative descriptors at low dimensions (512-dim), which is critical for storing large reference maps on the Jetson Nano.
\end{itemize}

\section{Implementation Flows for Current Architecture}

\subsection{Integration in the Benchmark Path}
To evaluate a new extractor using \code{demo.py} or \code{test_campus_dataset.py}, follow this flow:
\begin{enumerate}
    \item \textbf{Add Module:} Create a new wrapper in \code{feature_extraction/}.
    \item \textbf{Unified Interface:} Ensure the module follows the \code{compute_features(imgs)} contract used in \code{demo.py}.
    \item \textbf{Comparison:} Run the benchmark against the \code{GThard} and \code{GTsoft} matrices to generate Precision-Recall (PR) curves.
\end{enumerate}

\subsection{Integration in the Live Path}
The Live path building and retrieval logic must be updated across three modules:

\subsubsection{Phase 1: Factory Implementation}
Modify \code{live_vpr/extractors.py} to include a factory that initializes the model:
\begin{lstlisting}[style=pythoncode]
def get_extractor(model_name, config):
    if model_name == "selavpr_plus":
        return SelaVPRPlusExtractor(checkpoint=config.path)
    elif model_name == "vprtempo":
        return VPRTempoExtractor()
    #... handle other extractors
\end{lstlisting}

\subsubsection{Phase 2: Offline Map Building (\code{offline.py})}
When running in \code{build_map} mode, the descriptor matrix stored in the \code{.npz} file must match the new dimensions:
\begin{enumerate}
    \item Load sampled images from \code{artifacts/sampled_frames/}.
    \item \textbf{SelaVPR++ Note:} If using SelaVPR++, store \textit{both} the 512-bit binary hashes and the 4096-dim float vectors in the map metadata.
\end{enumerate}

\subsubsection{Phase 3: Online Localization (\code{online.py} and \code{search.py})}
The \code{LiveLocalizer} in \code{online.py} delegates search to \code{search.py}.
\begin{enumerate}
    \item \textbf{Exact Search:} For MixVPR and AnyLoc, the \code{exact} backend is sufficient.
    \item \textbf{Hamming Search:} For SelaVPR++, update \code{live_vpr/search.py} to implement a \code{HammingDistanceBackend} for the binary stage.
    \item \textbf{Re-ranking:} Use the \code{_rerank_exact} function in \code{search.py} to refine the top-K binary matches using the full float descriptors.
\end{enumerate}

\section{Suitability Matrix}
\begin{table}[H]
\centering
\begin{tabular}{lll}
\toprule
\textbf{Model} & \textbf{Deployment Target} & \textbf{Primary Solution} \\
\midrule
SelaVPR++ & Cloud-Edge Hybrid & Bandwidth efficiency via binary hashing \\
VPRTempo & On-Robot (CPU) & Real-time speed without GPU usage \\
SALAD & Cloud (High Precision) & Discarding dynamic campus noise \\
AnyLoc & Benchmark Baseline & Zero-shot transferability \\
MixVPR & On-Robot (GPU/RAM) & Low memory footprint for map storage \\
\bottomrule
\end{tabular}
\caption{Model recommendation based on campus robotic constraints.}
\end{table}

\end{document}